%*************************************
% ภาคผนวก ง.
%*************************************
\appendixChapterTitle{ง}{กระบวนการทำงาน NETCONF/YANG ในระบบ}
\vspace*{-0.8in}

\section{ภาพรวมโปรโตคอล NETCONF}

\begin{figure}[H]
\centering
\adjustbox{frame, width=1.0\textwidth}{
    \includegraphics[width=1.0\textwidth]{Image/netconf-1_2.png}
}
\caption{\fontSixTeen{ภาพรวมโปรโตคอล NETCONF}}
\label{figB:netconf1-2}
\end{figure}

\hspace*{1.3em}
NETCONF (Network Configuration Protocol) เปนโปรโตคอลมาตรฐานตาม RFC 6241 ทใชในการจดการคอนฟกอปกรณเครอขายผานชองทาง SSH ทปลอดภย โดยใช XML เปนรปแบบขอมลสำหรบการสอสารระหวาง Client (ระบบเวบแอปพลเคชน) กบ Server (อปกรณเครอขาย) ในโครงงานนระบบใช NETCONF เปนกลไกหลกในการเชอมตอและสงงานอปกรณ Cisco Nexus (NX-OS) และ IOS-XE

\hspace*{1.3em}
โปรโตคอล NETCONF มสถาปตยกรรมแบบ 4 ชน (Layer) ดงน

\begin{mycustomenum2}
    \item \textbf{Transport Layer:} ใช SSH (Secure Shell) เปนชองทางการสอสาร โดยเชอมตอผาน SSH Subsystem ชอ \textbf{netconf} ทพอรต 830 (คาเรมตนมาตรฐาน)
    \item \textbf{Messages Layer:} ใช XML-based RPC (Remote Procedure Call) ในการสงคำสงและรบผลลพธ โดยขอความทกขอความสนสดดวย Delimiter \textbf{]]>]]>} เพอระบจดสนสดของขอความ
    \item \textbf{Operations Layer:} กำหนดชดคำสง (Operations) ทสามารถดำเนนการได เชน \textbf{get-config}, \textbf{edit-config}, \textbf{lock}, \textbf{unlock}, \textbf{commit}, \textbf{validate} และ \textbf{close-session}
    \item \textbf{Content Layer:} ขอมลคอนฟกทจดการผาน NETCONF จะถกกำหนดโครงสรางดวย YANG Data Model
\end{mycustomenum2}

\section{YANG Data Model}
\hspace*{1.3em}
YANG (Yet Another Next Generation) เปนภาษามาตรฐานตาม RFC 7950 สำหรบกำหนดโครงสรางขอมล (Data Model) ของอปกรณเครอขาย YANG กำหนดวาขอมลคอนฟกทจะสงผาน NETCONF ตองมรปแบบอยางไร ประกอบดวยฟลดอะไรบาง แตละฟลดมชนดขอมลอะไร และมขอจำกดอยางไร

\hspace*{1.3em}
ในโครงงานนรองรบ YANG Namespace ของอปกรณ Cisco 2 ประเภทหลก ดงน

\begin{table}[H]
\centering
\caption{\fontSixTeen{YANG Namespace และ Root Element สำหรบอปกรณ Cisco ทรองรบในระบบ}}
\label{tabD:YangNamespaces}
\begin{tabularx}{\textwidth}{|l|X|l|}
\hline
\textbf{ประเภทอปกรณ} & \textbf{YANG Namespace} & \textbf{Root Element} \\
\hline
Cisco NX-OS (Nexus) & \textbf{http://cisco.com/ns/yang/cisco-nx-os-device} & \textbf{System} \\
\hline
Cisco IOS-XE & \textbf{http://cisco.com/ns/yang/Cisco-IOS-XE-native} & \textbf{native} \\
\hline
\end{tabularx}
\end{table}

\hspace*{1.3em}
โครงสราง YANG Model ประกอบดวยองคประกอบหลกดงน

\begin{mycustomenum2}
    \item \textbf{module:} หนวยหลกของ YANG Model กำหนดชอโมดลและ Namespace
    \item \textbf{namespace:} URI ทระบตวตนของโมดล ใชใน XML เพอระบวาขอมลเปนของโมดลใด
    \item \textbf{prefix:} ชอยอสำหรบอางอง Namespace ในเอกสาร XML
    \item \textbf{container:} โหนดทใชจดกลมขอมลยอย (เปรยบเสมอนโฟลเดอร)
    \item \textbf{list:} โหนดทมขอมลหลายรายการ โดยแตละรายการใช \textbf{key} ในการระบตวตน
    \item \textbf{leaf:} โหนดปลายสดทเกบคาขอมลจรง เชน ชออนเทอรเฟซ, IP Address
    \item \textbf{leaf-list:} โหนดปลายสดทเกบคาไดหลายคา (เชน รายการ VLAN ทอนญาต)
\end{mycustomenum2}

\hspace*{1.3em}
ในระบบนผใชสามารถอปโหลด YANG Model เองได โดยระบบจะ Parse ขอมลจาก YANG Content อตโนมต ดงคา \textbf{module}, \textbf{namespace}, \textbf{prefix} และ \textbf{revision} มาจดเกบในฐานขอมล นอกจากนแตละ YANG Model สามารถม XML Templates เพอใชเปน Template สำหรบ LLM ในการสราง Configuration ได

\section{กระบวนการเชอมตอ NETCONF Session}
\hspace*{1.3em}
การเชอมตอ NETCONF Session ในระบบมขนตอนดงตอไปน

\subsection{ขนตอนท 1: SSH Connection}
\hspace*{1.3em}
ระบบเรมตนดวยการสราง SSH Connection ไปยงอปกรณเครอขายทพอรต NETCONF (คาเรมตนคอ 830) โดยใช SSH2 Library ใน Node.js พรอมกำหนด Algorithms ทเขากนไดกบอปกรณ Cisco ดงน

\begin{mycustomenum2}
    \item \textbf{Key Exchange:} \textbf{diffie-hellman-group14-sha256}, \textbf{ecdh-sha2-nistp256}
    \item \textbf{Cipher:} \textbf{aes256-ctr}, \textbf{aes128-ctr}
    \item \textbf{Connection Timeout:} 30 วนาท
\end{mycustomenum2}

\subsection{ขนตอนท 2: NETCONF Subsystem}
\hspace*{1.3em}
เมอ SSH Connection สำเรจ ระบบจะเปด NETCONF Subsystem ผานคำสง \textbf{conn.subsys('netconf')} เพอเขาสโหมดการสอสาร NETCONF โดยเฉพาะ ซงจะได Stream สำหรบรบ-สงขอความ XML

\subsection{ขนตอนท 3: Hello Exchange}
\hspace*{1.3em}
ขนตอนแรกของโปรโตคอล NETCONF คอการแลกเปลยน Hello Message ระหวาง Client กบ Server เพอตกลง Capabilities ททงสองฝายรองรบ

\hspace*{1.3em}
\textbf{Server Hello:} อปกรณเครอขายจะสง Hello Message มากอน ระบบจะ Parse เพอดง Capabilities ทอปกรณรองรบ เชน \textbf{base:1.0}, \textbf{writable-running}, \textbf{candidate} เปนตน

\hspace*{1.3em}
\textbf{Client Hello:} ระบบตอบกลบดวย Hello Message ทประกาศ Capabilities ท Client รองรบ ตวอยาง Hello Message ทระบบสงไปมดงน

{\fontsize{13}{14}\selectfont
\begin{verbatim}
<?xml version="1.0" encoding="UTF-8"?>
<hello xmlns="urn:ietf:params:xml:ns:netconf:base:1.0">
  <capabilities>
    <capability>urn:ietf:params:netconf:base:1.0</capability>
    <capability>urn:ietf:params:netconf:capability:writable-running:1.0</capability>
    <capability>urn:ietf:params:netconf:capability:candidate:1.0</capability>
    <!-- ... other capabilities ... -->
  </capabilities>
</hello>]]>]]>
\end{verbatim}
}

\hspace*{1.3em}
เมอทงสองฝายแลกเปลยน Hello Message สำเรจ NETCONF Session จะถกสรางขนและพรอมสำหรบการดำเนนการ RPC ระบบจะจดเกบ Session ไวใน \textbf{connections Map} พรอมขอมล Capabilities, Stream และ Timestamp เพอใชในการดำเนนการตอไป

\section{NETCONF RPC Operations ทใชในระบบ}
\hspace*{1.3em}
ระบบใช NETCONF RPC Operations หลายชนดในการจดการคอนฟกอปกรณเครอขาย ทก RPC Message จะถกหอดวย \textbf{<rpc>} Envelope พรอมระบ \textbf{message-id} ทไมซำกน และสนสดดวย Delimiter \textbf{]]>]]>} โดยรปแบบทวไปของ RPC Message มดงน

{\fontsize{13}{14}\selectfont
\begin{verbatim}
<?xml version="1.0" encoding="UTF-8"?>
<rpc message-id="msg-1" xmlns="urn:ietf:params:xml:ns:netconf:base:1.0">
  <!-- Operation content here -->
</rpc>]]>]]>
\end{verbatim}
}

\hspace*{1.3em}
รายละเอยดของ RPC Operations แตละชนดมดงน

\subsection{get-config (อานคาคอนฟก)}
\hspace*{1.3em}
ใชสำหรบอานคา Running Configuration จากอปกรณ สามารถเลอกอานทงหมดหรอเฉพาะสวนทตองการดวย Subtree Filter ตวอยางการอานคาคอนฟกทงหมดจาก Running Datastore มดงน

{\fontsize{13}{14}\selectfont
\begin{verbatim}
<rpc message-id="1" xmlns="urn:ietf:params:xml:ns:netconf:base:1.0">
  <get-config><source><running/></source></get-config>
</rpc>]]>]]>
\end{verbatim}
}

\hspace*{1.3em}
ตวอยางการอานเฉพาะสวน Interface ของอปกรณ NX-OS ดวย Subtree Filter มดงน

{\fontsize{13}{14}\selectfont
\begin{verbatim}
<rpc message-id="2" xmlns="urn:ietf:params:xml:ns:netconf:base:1.0">
  <get-config>
    <source><running/></source>
    <filter type="subtree">
      <System xmlns="http://cisco.com/ns/yang/cisco-nx-os-device">
        <intf-items/>
      </System>
    </filter>
  </get-config>
</rpc>]]>]]>
\end{verbatim}
}

\subsection{edit-config (แกไขคอนฟก)}
\hspace*{1.3em}
ใชสำหรบแกไขคา Configuration ของอปกรณ สามารถเลอก Target Datastore เปน \textbf{running} (แกไขทนท) หรอ \textbf{candidate} (แกไขเปนรางกอน) และกำหนด Default Operation เปน \textbf{merge} (รวมกบคาเดม) หรอ \textbf{replace} (แทนทคาเดม)

\hspace*{1.3em}
ตวอยางการตงคา Hostname ของอปกรณ NX-OS ผาน edit-config มดงน

{\fontsize{13}{14}\selectfont
\begin{verbatim}
<rpc message-id="3" xmlns="urn:ietf:params:xml:ns:netconf:base:1.0">
  <edit-config>
    <target><running/></target>
    <default-operation>merge</default-operation>
    <config>
      <System xmlns="http://cisco.com/ns/yang/cisco-nx-os-device">
        <name>Switch-Core-01</name>
      </System>
    </config>
  </edit-config>
</rpc>]]>]]>
\end{verbatim}
}

\hspace*{1.3em}
ตวอยางการตงคา Interface ของอปกรณ IOS-XE ผาน edit-config มดงน

{\fontsize{13}{14}\selectfont
\begin{verbatim}
<rpc message-id="4" xmlns="urn:ietf:params:xml:ns:netconf:base:1.0">
  <edit-config>
    <target><running/></target>
    <default-operation>merge</default-operation>
    <config>
      <native xmlns="http://cisco.com/ns/yang/Cisco-IOS-XE-native">
        <interface>
          <GigabitEthernet>
            <name>1</name>
            <ip>
              <address>
                <primary>
                  <address>192.168.1.1</address>
                  <mask>255.255.255.0</mask>
                </primary>
              </address>
            </ip>
          </GigabitEthernet>
        </interface>
      </native>
    </config>
  </edit-config>
</rpc>]]>]]>
\end{verbatim}
}

\subsection{lock / unlock (ลอกและปลดลอก Datastore)}
\hspace*{1.3em}
ใชสำหรบลอก Datastore เพอปองกนไมให Session อนแกไขคาในขณะทกำลงทำงาน และปลดลอกเมอดำเนนการเสรจสน

{\fontsize{13}{14}\selectfont
\begin{verbatim}
<!-- Lock -->
<rpc message-id="5" xmlns="urn:ietf:params:xml:ns:netconf:base:1.0">
  <lock><target><candidate/></target></lock>
</rpc>]]>]]>

<!-- Unlock -->
<rpc message-id="6" xmlns="urn:ietf:params:xml:ns:netconf:base:1.0">
  <unlock><target><candidate/></target></unlock>
</rpc>]]>]]>
\end{verbatim}
}

\subsection{commit (ยนยนการเปลยนแปลง)}
\hspace*{1.3em}
ใชสำหรบยนยน (Commit) การเปลยนแปลงจาก Candidate Datastore ไปยง Running Datastore ทำใหคอนฟกมผลบงคบใชจรง

{\fontsize{13}{14}\selectfont
\begin{verbatim}
<rpc message-id="7" xmlns="urn:ietf:params:xml:ns:netconf:base:1.0">
  <commit/>
</rpc>]]>]]>
\end{verbatim}
}

\subsection{validate (ตรวจสอบความถกตอง)}
\hspace*{1.3em}
ใชตรวจสอบวาคอนฟกทเตรยมไวใน Candidate Datastore มความถกตองตาม YANG Schema หรอไม กอนทจะทำการ Commit

{\fontsize{13}{14}\selectfont
\begin{verbatim}
<rpc message-id="8" xmlns="urn:ietf:params:xml:ns:netconf:base:1.0">
  <validate><source><candidate/></source></validate>
</rpc>]]>]]>
\end{verbatim}
}

\subsection{discard-changes (ยกเลกการเปลยนแปลง)}
\hspace*{1.3em}
ใชสำหรบยกเลกการเปลยนแปลงทอยใน Candidate Datastore กลบไปเปนคาเดมของ Running Datastore ใชในกรณทตรวจพบขอผดพลาดกอนทจะ Commit

{\fontsize{13}{14}\selectfont
\begin{verbatim}
<rpc message-id="9" xmlns="urn:ietf:params:xml:ns:netconf:base:1.0">
  <discard-changes/>
</rpc>]]>]]>
\end{verbatim}
}

\subsection{close-session (ปด NETCONF Session)}
\hspace*{1.3em}
ใชสำหรบปด NETCONF Session อยางสมบรณเมอดำเนนการเสรจสน

{\fontsize{13}{14}\selectfont
\begin{verbatim}
<rpc message-id="10" xmlns="urn:ietf:params:xml:ns:netconf:base:1.0">
  <close-session/>
</rpc>]]>]]>
\end{verbatim}
}

\hspace*{1.3em}
ตารางท \ref{tabD:RpcOperations} สรป RPC Operations ทงหมดทใชในระบบพรอมคา Timeout ทกำหนด

\begin{table}[H]
\centering
\caption{\fontSixTeen{สรป NETCONF RPC Operations ทใชในระบบ}}
\label{tabD:RpcOperations}
\begin{tabularx}{\textwidth}{|l|l|l|X|}
\hline
\textbf{Operation} & \textbf{RPC} & \textbf{Timeout} & \textbf{คำอธบาย} \\
\hline
เชอมตอ & Hello Exchange & 30 วนาท & สราง SSH + NETCONF Session \\
\hline
อานคอนฟก & \textbf{get-config} & 180 วนาท & อาน Running Configuration \\
\hline
อานขอมล & \textbf{get} & 120 วนาท & อาน Operational Data \\
\hline
แกไขคอนฟก & \textbf{edit-config} & 90 วนาท & แกไข Running/Candidate \\
\hline
ยนยน & \textbf{commit} & 120 วนาท & ยนยน Candidate $\rightarrow$ Running \\
\hline
ตรวจสอบ & \textbf{validate} & 60 วนาท & ตรวจสอบ Candidate \\
\hline
ยกเลก & \textbf{discard-changes} & 15 วนาท & ยกเลก Candidate \\
\hline
ลอก/ปลดลอก & \textbf{lock/unlock} & 30 วนาท & ลอก Datastore \\
\hline
ปด Session & \textbf{close-session} & 10 วนาท & ปด NETCONF Session \\
\hline
\end{tabularx}
\end{table}

\section{กระบวนการ Deploy Configuration ผาน NETCONF}
\hspace*{1.3em}
การ Deploy Configuration ผาน NETCONF ในระบบม 2 รปแบบ ขนอยกบ Capabilities ทอปกรณรองรบ ดงน

\subsection{รปแบบท 1: ใช Candidate Datastore (แนะนำ)}
\hspace*{1.3em}
วธนเปนวธทปลอดภยทสดตามมาตรฐาน RFC 6241 เนองจากการเปลยนแปลงจะถกเขยนไปยง Candidate Datastore กอน และจะมผลบงคบใชกตอเมอทำการ Commit สำเรจเทานน หากเกดขอผดพลาดสามารถ Discard Changes เพอยอนกลบได ขนตอนมดงน

\begin{mycustomenum2}
    \item \textbf{Lock Candidate:} ลอก Candidate Datastore เพอปองกน Session อนแกไข
    \item \textbf{Edit-Config:} สง Configuration XML ไปยง Candidate Datastore ดวย \textbf{edit-config}
    \item \textbf{Validate (ถารองรบ):} ตรวจสอบความถกตองของ Candidate ดวย \textbf{validate}
    \item \textbf{Commit:} ยนยนการเปลยนแปลงจาก Candidate ไปยง Running ดวย \textbf{commit}
    \item \textbf{Unlock Candidate:} ปลดลอก Candidate Datastore
\end{mycustomenum2}

\hspace*{1.3em}
หากขนตอนใดเกดขอผดพลาด ระบบจะทำ Cleanup อตโนมต ไดแก Discard Changes เพอยกเลกการเปลยนแปลงใน Candidate และ Unlock เพอปลดลอก Datastore

\subsection{รปแบบท 2: ใช Writable-Running (Fallback)}
\hspace*{1.3em}
สำหรบอปกรณทไมรองรบ Candidate Datastore ระบบจะใช \textbf{writable-running} เปนทางเลอก โดยการเปลยนแปลงจะมผลทนทเมอ edit-config สำเรจ ขนตอนมดงน

\begin{mycustomenum2}
    \item \textbf{Edit-Config:} สง Configuration XML ไปยง Running Datastore โดยตรง
    \item การเปลยนแปลงมผลบงคบใชทนท ไมตอง Commit
\end{mycustomenum2}

\hspace*{1.3em}
\textbf{ขอควรระวง:} วธนไมมขนตอน Validate กอน Commit และไมสามารถ Rollback ไดหากเกดขอผดพลาด จงควรใช Candidate Datastore เปนลำดบแรกเมอเปนไปได

\section{Subtree Filter สำหรบอานขอมลอปกรณ}
\hspace*{1.3em}
ระบบใช Subtree Filter ใน \textbf{<get>} RPC เพออานขอมลเฉพาะสวนทตองการจากอปกรณ โดยไมตองอานขอมลทงหมด ซงชวยลดปรมาณขอมลทรบ-สงและเพมประสทธภาพ ระบบรองรบ Filter หลายประเภทตามชนดของอปกรณ ดงน

\subsection{Subtree Filter สำหรบ NX-OS (Nexus)}
\hspace*{1.3em}
YANG Model ของ Cisco NX-OS ใช Root Element \textbf{System} ภายใต Namespace \textbf{cisco-nx-os-device} ตวอยาง Filter ทใชในระบบมดงน

\begin{table}[H]
\centering
\caption{\fontSixTeen{Subtree Filter สำหรบอปกรณ Cisco NX-OS}}
\label{tabD:NxosFilters}
\begin{tabularx}{\textwidth}{|l|l|X|}
\hline
\textbf{ประเภทขอมล} & \textbf{YANG Path} & \textbf{คำอธบาย} \\
\hline
System Info & \textbf{System/name} & ชออปกรณ (Hostname) \\
\hline
Interfaces & \textbf{System/intf-items} & อนเทอรเฟซทงหมด (Physical, SVI, Port-channel) \\
\hline
VLANs & \textbf{System/bd-items} & Bridge Domain / VLAN \\
\hline
OSPF & \textbf{System/ospf-items} & Routing Protocol OSPF \\
\hline
BGP & \textbf{System/bgp-items} & Routing Protocol BGP \\
\hline
Features & \textbf{System/fm-items} & Feature Manager (เปด/ปดฟเจอร) \\
\hline
\end{tabularx}
\end{table}

\subsection{Subtree Filter สำหรบ IOS-XE}
\hspace*{1.3em}
YANG Model ของ Cisco IOS-XE ใช Root Element \textbf{native} ภายใต Namespace \textbf{Cisco-IOS-XE-native} ตวอยาง Filter ทใชในระบบมดงน

\begin{table}[H]
\centering
\caption{\fontSixTeen{Subtree Filter สำหรบอปกรณ Cisco IOS-XE}}
\label{tabD:IosXeFilters}
\begin{tabularx}{\textwidth}{|l|l|X|}
\hline
\textbf{ประเภทขอมล} & \textbf{YANG Path} & \textbf{คำอธบาย} \\
\hline
Hostname & \textbf{native/hostname} & ชออปกรณ \\
\hline
Interfaces & \textbf{native/interface} & อนเทอรเฟซ GigabitEthernet/Loopback \\
\hline
VLANs & \textbf{native/vlan} & VLAN Configuration \\
\hline
OSPF & \textbf{native/router/ospf} & Routing Protocol OSPF \\
\hline
BGP & \textbf{native/router/bgp} & Routing Protocol BGP \\
\hline
\end{tabularx}
\end{table}

\section{การ Deploy ผาน Agent Relay}
\hspace*{1.3em}
เนองจากระบบ Backend ทำงานเปน Vercel Serverless Functions ซงไมสามารถเชอมตอ SSH/NETCONF ไปยงอปกรณในเครอขายภายในไดโดยตรง ระบบจงใช Desktop Agent เปนตวกลาง (Relay) ในการสงคำสง NETCONF ไปยงอปกรณ กระบวนการมดงน

\begin{mycustomenum2}
    \item \textbf{Backend สราง Command:} เมอผใชสง Deploy ระบบจะสราง AgentCommand ใน MongoDB พรอมขอมล NETCONF XML Configuration
    \item \textbf{Agent Poll Command:} Desktop Agent ทำ HTTP Polling ทก 2 วนาทเพอดง Pending Commands
    \item \textbf{Agent ดำเนนการ:} Agent เชอมตอ SSH/NETCONF ไปยงอปกรณจรงในเครอขายภายใน ดำเนนการตาม Workflow (Lock $\rightarrow$ Edit-Config $\rightarrow$ Validate $\rightarrow$ Commit $\rightarrow$ Unlock)
    \item \textbf{Agent รายงานผล:} Agent สงผลลพธกลบมายง Backend ผาน HTTP POST
    \item \textbf{Backend สงผลใหผใช:} Backend อปเดตสถานะ Command และแจงผลลพธใหผใชทราบ
\end{mycustomenum2}

\section{การจดการ YANG Model ในระบบ}
\hspace*{1.3em}
ระบบมระบบจดการ YANG Model ทชวยใหผใชสามารถอปโหลด จดเกบ และใช YANG Model ในการสราง Configuration ผาน LLM ได โดยมองคประกอบหลกดงน

\subsection{การอปโหลดและ Parse YANG Model}
\hspace*{1.3em}
เมอผใชอปโหลด YANG Model ระบบจะ Parse เนอหาอตโนมตเพอดงขอมลสำคญ ไดแก

\begin{mycustomenum2}
    \item \textbf{module name:} ชอโมดล YANG เชน \textbf{cisco-nx-os-device}
    \item \textbf{namespace:} URI ของโมดล เชน \textbf{http://cisco.com/ns/yang/cisco-nx-os-device}
    \item \textbf{prefix:} ชอยอสำหรบอางอง เชน \textbf{nxos}
    \item \textbf{revision:} เวอรชนของ YANG Model เชน \textbf{2024-01-01}
\end{mycustomenum2}

\subsection{XML Templates}
\hspace*{1.3em}
แตละ YANG Model สามารถม XML Templates หลายรายการ ซงเปนตวอยาง XML Configuration ทถกตองตาม YANG Schema ระบบจะสง Templates เหลานให LLM เปนตวอยางประกอบเพอชวยให LLM สราง Configuration ทม Syntax และ Namespace ถกตอง

\subsection{การใช YANG Model สำหรบ LLM Generation}
\hspace*{1.3em}
เมอผใชสราง Configuration ดวย LLM ระบบจะ

\begin{mycustomenum2}
    \item ดง YANG Models ทตรงกบประเภทอปกรณ (NX-OS หรอ IOS-XE) จากฐานขอมล
    \item จดรปแบบขอมล YANG Model เปนขอความบรบท (Context) ประกอบดวย Namespace, Config Paths และ XML Templates
    \item สงขอมลบรบทพรอม Prompt ของผใชไปให LLM เพอสราง Configuration ทถกตองตาม YANG Schema
\end{mycustomenum2}

\section{การจดการ Error ใน NETCONF}
\hspace*{1.3em}
เมออปกรณพบขอผดพลาดในการดำเนนการ NETCONF จะตอบกลบดวย \textbf{<rpc-error>} ระบบจะ Parse ขอมลใน Error Message เพอแสดงรายละเอยดใหผใชทราบ โดย Error Message ประกอบดวย

\begin{mycustomenum2}
    \item \textbf{error-tag:} ประเภทของขอผดพลาด เชน \textbf{invalid-value}, \textbf{data-missing}, \textbf{operation-failed}
    \item \textbf{error-message:} คำอธบายขอผดพลาดเปนขอความภาษาองกฤษ
    \item \textbf{error-path:} ตำแหนงใน YANG Model ทเกดขอผดพลาด เชน \textbf{/System/intf-items}
\end{mycustomenum2}

\hspace*{1.3em}
ตวอยาง RPC Error Response มดงน

{\fontsize{13}{14}\selectfont
\begin{verbatim}
<rpc-reply xmlns="urn:ietf:params:xml:ns:netconf:base:1.0" message-id="3">
  <rpc-error>
    <error-type>application</error-type>
    <error-tag>invalid-value</error-tag>
    <error-severity>error</error-severity>
    <error-path>/System/intf-items</error-path>
    <error-message xml:lang="en">Invalid interface name</error-message>
  </rpc-error>
</rpc-reply>]]>]]>
\end{verbatim}
}

\hspace*{1.3em}
ในกระบวนการ Deploy ผาน Candidate Datastore หากเกดขอผดพลาดตรงขนตอนใดกตาม ระบบจะทำ Cleanup อตโนมต ประกอบดวย Discard Changes เพอยกเลกการเปลยนแปลงใน Candidate, Unlock เพอปลดลอก Datastore และแจง Error กลบใหผใชทราบพรอมรายละเอยด เพอใหสามารถแกไข Configuration และลองใหมได
